\documentclass[12pt]{article}
\usepackage[utf8]{inputenc}
\usepackage[T1]{fontenc}
\usepackage{hyperref}
\title{Vute Language Resources for K-12 Curriculum Development}
\author{Compiled by Claude Sonnet 5 (claude-sonnet-5) for Andr\'as Kornai}
\date{}
\begin{document}
\maketitle
\section{Introduction}
Vute (Wute, Voute; ISO 639-3 vut; Glottolog vute1244) is a Mambiloid (Northern Bantoid) language with about 21,000 speakers, mostly in the Mayo-Banyo and Dj\'erem departments of Cameroon's Adamawa Region, and a few in southeastern Nigeria. Dialects are Eastern, Central and Doume. A Latin-script orthography was established in 1979.
\section{Primary Educational Resources}
No primers or school textbooks were found in search results.
\section{Linguistic and Reference Materials}
An overview of the Vute writing system is at \url{https://www.omniglot.com/writing/vute.htm}. A lexical resource page is at \url{https://cocoon.huma-num.fr/exist/crdo/lexvo.xql?code=vut}.
\section{Religious and Community Resources}
A New Testament (2007, Bible Society of Cameroon and Lutheran Bible Translators) is reported. Scripture listing: \url{https://www.scriptureearth.org/00eng.php?iso=vut}. Translation notes: \url{https://tips.translation.bible/tip_language/vut/}. Joshua Project: \url{https://www.joshuaproject.net/index.php/people_groups/15807/CM}
\section{Digital and Audio Resources}
Global Recordings Network offers audio and video in Vute: \url{https://globalrecordings.net/en/language/vut}
\section{Curriculum Development Context}
An established orthography and a New Testament give a base for literacy materials, but no school programme or textbooks were found. I did not verify a full Bible.
\end{document}